\documentclass[11pt,a4paper]{article}

\usepackage{amsmath,amsthm,amssymb}
\usepackage{booktabs}
\usepackage[ruled,vlined,linesnumbered]{algorithm2e}
\usepackage[colorlinks=true,linkcolor=blue,citecolor=blue,urlcolor=blue]{hyperref}
\usepackage{geometry}
\geometry{margin=2.5cm}
\usepackage{microtype}

% ---------- theorem environments ----------
\newtheorem{theorem}{Theorem}[section]
\newtheorem{proposition}[theorem]{Proposition}
\newtheorem{corollary}[theorem]{Corollary}
\newtheorem{definition}[theorem]{Definition}
\newtheorem{remark}[theorem]{Remark}

% ---------- notation macros ----------
\newcommand{\R}{\mathbb{R}}
\newcommand{\bA}{\mathbf{A}}
\newcommand{\bS}{\mathbf{S}}
\newcommand{\bb}{\mathbf{b}}
\newcommand{\bp}{\mathbf{p}}
\newcommand{\bg}{\mathbf{g}}
\newcommand{\br}{\mathbf{r}}
\newcommand{\bd}{\mathbf{d}}
\newcommand{\bt}{\mathbf{t}}
\newcommand{\bu}{\mathbf{u}}
\newcommand{\bw}{\mathbf{w}}
\newcommand{\bx}{\mathbf{x}}
\newcommand{\by}{\mathbf{y}}
\newcommand{\bZ}{\mathbf{Z}}
\newcommand{\be}{\mathbf{e}}
\newcommand{\bone}{\mathbf{1}}
\newcommand{\bzero}{\mathbf{0}}
\newcommand{\Proj}{\mathcal{P}}
\newcommand{\calC}{\mathcal{C}}
\newcommand{\calF}{\mathcal{F}}
\newcommand{\calA}{\mathcal{A}}
\newcommand{\norm}[1]{\|#1\|}
\newcommand{\inner}[2]{\langle #1,\, #2 \rangle}
\newcommand{\tr}{\top}

% ---------- algorithm2e style ----------
\SetAlgoLined
\DontPrintSemicolon
\SetKwComment{Comment}{// }{}

% =========================================================
\begin{document}

\title{\textbf{Projected Conjugate Gradient Methods\\
       for Bound-Constrained Quadratic Programs}\\[0.5em]
       \large Theory, Line Search Strategies, and Application\\
       to Elastic Contact Mechanics}
\author{Hcontact project --- technical reference}
\date{June 2026}
\maketitle
\tableofcontents

% ---------------------------------------------------------
\section{The Conjugate Gradient Method for Quadratic Minimisation}
\label{sec:cg}

\subsection{Problem statement}

Consider the unconstrained quadratic program
\begin{equation}
  \min_{\bx \in \R^n}\; f(\bx) = \tfrac{1}{2}\bx^\tr \bA \bx - \bb^\tr \bx,
  \label{eq:qp}
\end{equation}
where $\bA \in \R^{n \times n}$ is symmetric positive definite (SPD).
The unique minimiser $\bx^*$ satisfies $\nabla f(\bx^*) = \bA\bx^* - \bb = \bzero$,
i.e.\ it is the solution of the linear system $\bA\bx = \bb$.

\subsection{Algorithm}

The conjugate gradient (CG) method of Hestenes and Stiefel~\cite{HestStief1952}
generates iterates $\bx_0, \bx_1, \ldots$ as follows.

\begin{definition}[Residual and search direction]
  Given an iterate $\bx_k$, define the \emph{residual} $\br_k = \bb - \bA\bx_k = -\nabla f(\bx_k)$
  and the \emph{search direction} $\bd_k \in \R^n$.
\end{definition}

Starting from $\bx_0$ (often $\bx_0 = \bzero$), with $\br_0 = \bb - \bA\bx_0$ and
$\bd_0 = \br_0$, each iteration computes:
\begin{align}
  \alpha_k   &= \frac{\br_k^\tr \br_k}{\bd_k^\tr \bA \bd_k},
              \label{eq:cg-alpha}\\[4pt]
  \bx_{k+1} &= \bx_k + \alpha_k \bd_k,
              \label{eq:cg-x}\\[4pt]
  \br_{k+1} &= \br_k - \alpha_k \bA \bd_k,
              \label{eq:cg-r}\\[4pt]
  \beta_k    &= \frac{\br_{k+1}^\tr \br_{k+1}}{\br_k^\tr \br_k},
              \label{eq:cg-beta-FR}\\[4pt]
  \bd_{k+1} &= \br_{k+1} + \beta_k \bd_k.
              \label{eq:cg-d}
\end{align}

The step length $\alpha_k$ in~\eqref{eq:cg-alpha} is the \emph{exact minimiser}
of $\phi(\alpha) = f(\bx_k + \alpha\bd_k)$ along the ray, since
\[
  \phi'(\alpha) = 0 \;\Longrightarrow\; \alpha_k
  = \frac{-\nabla f(\bx_k)^\tr \bd_k}{\bd_k^\tr \bA \bd_k}
  = \frac{\br_k^\tr \bd_k}{\bd_k^\tr \bA \bd_k}.
\]
At $k=0$ this reduces to $\alpha_0 = \norm{\br_0}^2 / (\br_0^\tr \bA \br_0)$
since $\bd_0 = \br_0$ and $\br_0^\tr\bd_0 = \norm{\br_0}^2$.

\subsection{Key algebraic properties}

\begin{proposition}[Orthogonality and conjugacy]
  \label{prop:orthogonality}
  The residuals and search directions generated by CG satisfy:
  \begin{enumerate}
    \item $\br_i^\tr \br_j = 0$ for $i \neq j$ (mutual orthogonality of residuals).
    \item $\bd_i^\tr \bA \bd_j = 0$ for $i \neq j$ ($\bA$-conjugacy of directions).
    \item $\br_{k+1}^\tr \bd_k = 0$ (new residual orthogonal to last direction).
  \end{enumerate}
\end{proposition}

\begin{proof}
  By induction on $k$, using~\eqref{eq:cg-r} and~\eqref{eq:cg-d}; see~\cite{NocWri2006} for details.
\end{proof}

The recurrence $\br_{k+1} = \br_k - \alpha_k\bA\bd_k$ avoids storing all previous directions:
only the last residual and direction are needed.

\begin{theorem}[Finite termination]
  \label{thm:finite}
  In exact arithmetic, the CG iterates satisfy $f(\bx_k) \geq f(\bx_{k+1})$ for all $k$,
  and $\bx_k = \bx^*$ for some $k \leq n$.
  If $\bA$ has only $m \leq n$ distinct eigenvalues, termination occurs in at most $m$ steps.
\end{theorem}

\begin{remark}
  Each iteration costs one matrix-vector product $\bA\bd_k$ plus $O(n)$ arithmetic.
  For large systems solved via an H-matrix, the matvec dominates and has cost
  $O(n \log n)$ in storage and $O(n \log n)$ in flops.
\end{remark}

% ---------------------------------------------------------
\section{Beta Update Formulas}
\label{sec:beta}

\subsection{Motivation}

For a general (non-quadratic) objective, the exact formula~\eqref{eq:cg-beta-FR}
for $\beta_k$ no longer guarantees descent.  Moreover, in the bound-constrained
setting of Section~\ref{sec:polonsky}, the effective gradient changes discontinuously
whenever the active set (contact set) changes, making stale conjugation information
harmful.  This section surveys the principal $\beta$ formulas and their practical behaviour.

\subsection{Notation}

Let $\by_k = \br_{k+1} - \br_k$ denote the residual change.  In the quadratic case,
$\by_k = -\alpha_k\bA\bd_k$ and the orthogonality relations of
Proposition~\ref{prop:orthogonality} simplify all formulas to the same value.
The formulas differ for non-quadratic $f$ or for inexact line search.

\subsection{The four principal formulas}

\begin{table}[ht]
\centering
\caption{Standard nonlinear CG $\beta$ update formulas.}
\label{tab:beta}
\begin{tabular}{@{}lll@{}}
\toprule
Name & Formula for $\beta_k$ & Year \\
\midrule
Hestenes--Stiefel (HS)  & $\dfrac{\br_{k+1}^\tr \by_k}{\bd_k^\tr \by_k}$ & 1952 \\[8pt]
Fletcher--Reeves (FR)   & $\dfrac{\norm{\br_{k+1}}^2}{\norm{\br_k}^2}$ & 1964 \\[8pt]
Polak--Ribi\`ere (PR)   & $\dfrac{\br_{k+1}^\tr \by_k}{\norm{\br_k}^2}$ & 1969 \\[8pt]
Polak--Ribi\`ere+ (PR+) & $\max\!\bigl(0,\;\beta_k^{\mathrm{PR}}\bigr)$ & --- \\[8pt]
Dai--Yuan (DY)          & $\dfrac{\norm{\br_{k+1}}^2}{\bd_k^\tr \by_k}$ & 1999 \\
\bottomrule
\end{tabular}
\end{table}

\noindent
All formulas coincide in the quadratic case with exact line search:
$\by_k = -\alpha_k\bA\bd_k$, $\br_k^\tr\by_k = 0$ (by orthogonality), so
\[
  \beta_k^{\mathrm{PR}} = \frac{\norm{\br_{k+1}}^2 - \br_{k+1}^\tr\br_k}{\norm{\br_k}^2}
  = \frac{\norm{\br_{k+1}}^2}{\norm{\br_k}^2} = \beta_k^{\mathrm{FR}}.
\]

\subsection{Convergence and practical behaviour}

\paragraph{Fletcher--Reeves.}
FR is \emph{globally convergent} for any line search satisfying the strong
Wolfe conditions~\cite{AlBaali1985,Zoutendijk1970}.  However, when iterates
stagnate ($\norm{\br_{k+1}} \approx \norm{\br_k}$), FR keeps $\beta \approx 1$,
so the new direction barely changes and progress is slow.

\paragraph{Polak--Ribi\`ere and PR+.}
PR is \emph{not} globally convergent in general:
Gilbert and Nocedal~\cite{GilNoc1992} construct a counterexample.
Its practical advantage comes from \emph{automatic restarting}: whenever
$\br_{k+1} \approx \br_k$ (small progress), $\by_k \approx \bzero$ and
$\beta^{\mathrm{PR}} \approx 0$, so the next step is essentially steepest descent.
PR+ enforces this by clipping: $\beta^{\mathrm{PR+}} = \max(0, \beta^{\mathrm{PR}})$.

\begin{proposition}[Descent property of PR+]
  If the line search satisfies the strong Wolfe conditions then the PR+ direction
  satisfies $\nabla f(\bx_{k+1})^\tr \bd_{k+1} \leq -c\norm{\nabla f(\bx_{k+1})}^2$
  for some $c > 0$.  In particular it is a descent direction.
\end{proposition}

\paragraph{Hestenes--Stiefel.}
HS enforces $\bd_{k+1}^\tr \by_k = 0$ by construction, a form of conjugacy with
respect to the curvature $\by_k$.  It behaves similarly to PR for non-quadratic
problems, but can fail when $\bd_k^\tr\by_k \approx 0$.

\paragraph{Dai--Yuan.}
DY is always $\geq 0$ (since $\bd_k^\tr\by_k > 0$ under Wolfe conditions) and
globally convergent~\cite{DaiYuan1999}, but tends to be conservative (small $\beta$),
similar in behaviour to restarted steepest descent.

\paragraph{Summary.}
For the bound-constrained contact QP (Section~\ref{sec:polonsky}), the objective is
quadratic within a fixed contact set, so all $\beta$ formulas are equivalent there.
The distinction arises across active-set changes, which introduce a ``non-quadratic''
perturbation.  PR+ handles this gracefully: the automatic restart (via $\beta = 0$)
discards stale conjugation whenever the contact set changes, whereas FR retains
a stale $\beta \approx G/G_{\mathrm{old}}$ that may slow convergence.

% ---------------------------------------------------------
\section{Line Search Strategies}
\label{sec:linesearch}

Given a current iterate $\bx$ and a descent direction $\bd$
($\nabla f(\bx)^\tr\bd < 0$), a \emph{line search} selects a step length
$\alpha > 0$ such that $\bx + \alpha\bd$ reduces $f$ adequately.

Define $\phi(\alpha) = f(\bx + \alpha\bd)$ and
$\phi'(0) = \nabla f(\bx)^\tr\bd < 0$.

\subsection{Exact line search (quadratic objective)}

For the quadratic $f(\bx) = \tfrac{1}{2}\bx^\tr\bA\bx - \bb^\tr\bx$:
\begin{equation}
  \alpha^* = \operatorname{argmin}_{\alpha \geq 0}\,\phi(\alpha)
           = \frac{\br^\tr \bd}{\bd^\tr \bA \bd},
  \label{eq:exact-ls}
\end{equation}
where $\br = \bb - \bA\bx$.  This requires \emph{one matrix-vector product}
$\bA\bd$ (the same product needed to update the residual in~\eqref{eq:cg-r}).

\subsection{Armijo sufficient decrease}

Accept $\alpha$ if
\begin{equation}
  \phi(\alpha) \leq \phi(0) + c_1\,\alpha\,\phi'(0), \qquad c_1 \in (0,1).
  \label{eq:armijo}
\end{equation}
Typically $c_1 = 10^{-4}$.  The standard backtracking algorithm starts with
$\alpha = 1$ and repeatedly halves until~\eqref{eq:armijo} is satisfied.

\begin{proposition}
  If $\bA$ is SPD and $\bd$ is a descent direction, then the Armijo condition
  is satisfied for all $\alpha \in (0, \bar\alpha]$ for some $\bar\alpha > 0$
  determined by the Lipschitz constant of $\nabla f$.
\end{proposition}

\subsection{Wolfe and strong Wolfe conditions}

The \emph{Wolfe conditions} add a curvature requirement to Armijo:
\begin{equation}
  \phi'(\alpha) \geq c_2\,\phi'(0), \qquad c_1 < c_2 < 1.
  \label{eq:wolfe}
\end{equation}
The \emph{strong Wolfe conditions} replace~\eqref{eq:wolfe} by
$|\phi'(\alpha)| \leq c_2\,|\phi'(0)|$, preventing the step from being too
large (overshooting the minimum) as well as too small.  Typically $c_2 = 0.9$
for nonlinear CG.

The Wolfe conditions are essential for the convergence proofs of FR
(Al-Baali~\cite{AlBaali1985}) and DY~\cite{DaiYuan1999}.

\subsection{Polynomial interpolation}

Given $\phi(0)$, $\phi'(0)$, and a trial $\phi(\alpha_0)$ that fails
Armijo, fit a quadratic:
\[
  \hat\phi(\alpha) = a\alpha^2 + b\alpha + c, \quad
  c = \phi(0),\quad b = \phi'(0),\quad
  a = \frac{\phi(\alpha_0) - b\alpha_0 - c}{\alpha_0^2}.
\]
The minimiser is $\alpha^* = -b/(2a)$, clamped to $[0.01\,\alpha_0,\; 0.5\,\alpha_0]$
to avoid very small or very large steps.  This is the strategy used in the
GPCG gradient-projection phase~\cite{MoreToraldo1991} (Section~\ref{sec:gpcg}).

\subsection{Cost comparison}

\begin{table}[ht]
\centering
\caption{Approximate matvec cost per CG iteration for each line search strategy.}
\label{tab:ls-cost}
\begin{tabular}{@{}lll@{}}
\toprule
Strategy & Matvecs / iter & Notes \\
\midrule
Exact (quadratic)       & $2$ & 1 for gradient, 1 for $\bS\bt$; both unavoidable \\
Armijo backtrack        & $2$--$4$ & Each failed trial costs one extra evaluation \\
Wolfe bracket $+$ zoom  & $3$--$6$ & More robust for non-quadratic objectives \\
\bottomrule
\end{tabular}
\end{table}

The P-K algorithm (Section~\ref{sec:polonsky}) uses exact line search:
two matvecs per outer iteration are unavoidable for a quadratic objective
(one for $\nabla f$, one for the curvature along the search direction).

% ---------------------------------------------------------
\section{Projected Gradient and Active Sets}
\label{sec:projection}

We now turn to \emph{bound-constrained} QPs:
\begin{equation}
  \min_{\bx \in \R^n}\; f(\bx) = \tfrac{1}{2}\bx^\tr\bA\bx - \bb^\tr\bx
  \quad \text{s.t.}\quad l_i \leq x_i \leq u_i,\quad i = 1,\ldots,n.
  \label{eq:bqp}
\end{equation}

\subsection{Projection operator}

The orthogonal projection onto the feasible box $[l,u]^n$ is
\begin{equation}
  [\Proj(\bx)]_i = \max\!\bigl(l_i,\;\min(x_i, u_i)\bigr).
  \label{eq:projection}
\end{equation}

\subsection{Active and free sets}

\begin{definition}
  At a feasible point $\bx$, the \emph{active set}
  $\calA(\bx) = \{i : x_i = l_i \text{ or } x_i = u_i\}$ collects
  the indices at their bounds, and the \emph{free set}
  $\calF(\bx) = \{1,\ldots,n\} \setminus \calA(\bx)$ collects interior indices.
\end{definition}

\subsection{KKT optimality conditions}

$\bx^*$ is optimal for~\eqref{eq:bqp} if and only if there exist
multipliers $\lambda_i^l \geq 0$, $\lambda_i^u \geq 0$ such that
$\nabla f(\bx^*)_i = \lambda_i^l - \lambda_i^u$ and
$\lambda_i^l(x_i^* - l_i) = \lambda_i^u(u_i - x_i^*) = 0$,
or equivalently:
\begin{equation}
  [\nabla_P f(\bx^*)]_i = 0 \quad \forall i,
  \label{eq:kkt}
\end{equation}
where the \emph{projected gradient} is
\begin{equation}
  [\nabla_P f(\bx)]_i =
  \begin{cases}
    \min(\nabla f(\bx)_i,\; 0) & \text{if } x_i = l_i, \\
    \max(\nabla f(\bx)_i,\; 0) & \text{if } x_i = u_i, \\
    \nabla f(\bx)_i             & \text{otherwise.}
  \end{cases}
  \label{eq:proj-grad}
\end{equation}
The projected gradient combines the sign constraints from the KKT conditions
into a single vector, making it the natural optimality residual.

\subsection{Projected gradient step}

The simplest feasible descent step is
$\bx_{k+1} = \Proj[\bx_k - \alpha_k\nabla f(\bx_k)]$.
The projection simultaneously enforces feasibility and implicitly identifies
active constraints: any index $i$ for which $x_{k,i} - \alpha_k\nabla f(\bx_k)_i$
lies outside $[l_i, u_i]$ is ``clamped'' to its bound and enters $\calA(\bx_{k+1})$.

% ---------------------------------------------------------
\section{The Polonsky--Keer Projected CG Algorithm}
\label{sec:polonsky}

\subsection{The contact QP}

The Boussinesq elastic half-space contact problem with a uniform Cartesian grid
of $N = N_s^2$ nodes leads to the QP
\begin{equation}
  \min_{\bp \in \R^N}\; f(\bp) = \tfrac{1}{2}\bp^\tr\bS\bp + \bp^\tr\bg_0
  \quad \text{s.t.}\quad p_i \geq 0,\quad \frac{1}{N}\sum_{i=1}^N p_i = \bar p,
  \label{eq:contact-qp}
\end{equation}
where $\bS \in \R^{N\times N}$ is the SPD elastic influence matrix (Boussinesq
kernel, Love 1929 element formula), $\bg_0 \in \R^N$ is the initial surface gap,
and $\bar p > 0$ is the prescribed nominal contact pressure.
The gradient is $\nabla f(\bp) = \bS\bp + \bg_0 = \bu + \bg_0$,
where $\bu = \bS\bp$ is the elastic surface displacement.

\subsection{Handling the equality constraint: rigid-body shift}

Problem~\eqref{eq:contact-qp} couples a non-negativity constraint on $\bp$
with an equality constraint on its mean.
At the optimal solution, complementarity requires
\begin{equation}
  p_i > 0 \;\Rightarrow\; g_i = u_i + g_{0,i} - \delta = 0,
  \label{eq:complementarity}
\end{equation}
where $\delta$ is the \emph{rigid-body approach} (unknown).
Given the current iterate $\bp$, estimate $\delta$ from the contact set
$\calC = \{i : p_i > 0\}$:
\begin{equation}
  \delta = \frac{1}{|\calC|}\sum_{i \in \calC}(u_i + g_{0,i}),
  \qquad \tilde g_i = u_i + g_{0,i} - \delta.
  \label{eq:shift}
\end{equation}
The shifted gradient $\tilde{\bg}$ satisfies $\sum_{i\in\calC}\tilde g_i = 0$,
which enforces the mean-pressure constraint implicitly.
The equality $\sum_i p_i = N\bar p$ is enforced explicitly by a rescaling step
after each pressure update.

\subsection{CG direction restricted to the contact set}

The P-K algorithm restricts the CG search direction $\bt$ to the contact set:
$t_i = 0$ for $i \notin \calC$.  The $\beta$ update uses only the contact-set
components of $\tilde{\bg}$:
\begin{align}
  G_k &= \sum_{i \in \calC} \tilde g_i^2, \label{eq:Gk}\\[4pt]
  \beta_k^{\mathrm{FR}}  &= G_k / G_{k-1}, \label{eq:beta-fr-pk}\\[4pt]
  \beta_k^{\mathrm{PR+}} &= \max\!\Bigl(0,\;
     \tfrac{1}{G_{k-1}}\textstyle\sum_{i\in\calC}\tilde g_i(\tilde g_i - \tilde g_i^{\mathrm{prev}})\Bigr),
  \label{eq:beta-pr-pk}
\end{align}
where $\tilde{\bg}^{\mathrm{prev}}$ is the shifted gradient from the previous
iteration (over the current $\calC$).  When the contact set changes,
$\beta$ is reset to zero regardless of the formula chosen.

\subsection{Mean-corrected exact line search}

The exact line search must account for the rigid-body shift that will be applied
at the next iteration.  Given $\br = \bS\bt$ and $\bar r_\calC = |\calC|^{-1}\sum_{i\in\calC}r_i$:
\begin{equation}
  \tau = \frac{\sum_{i\in\calC}\tilde g_i\, t_i}
              {\sum_{i\in\calC}(r_i - \bar r_\calC)\, t_i}.
  \label{eq:tau}
\end{equation}
The denominator $r_i - \bar r_\calC$ removes the mean displacement from $\br$
on the contact set, matching the shift that will be subtracted from $\bg$ at the
start of the next iteration.

\subsection{Overlap correction}

After the projection step $p_i \leftarrow \max(p_i - \tau t_i, 0)$,
some nodes may drop to zero pressure while still having $\tilde g_i < 0$
(negative gap, i.e.\ physical penetration).
This contradicts the complementarity condition~\eqref{eq:complementarity}
and indicates that these nodes should be in the contact set.
The \emph{overlap correction}~\cite{PolKeer1999} restores them:
\begin{equation}
  p_i \leftarrow -\tau\tilde g_i \quad
  \text{for all } i \text{ with } p_i = 0 \text{ and } \tilde g_i < 0.
  \label{eq:overlap}
\end{equation}
Whenever~\eqref{eq:overlap} is applied, $\beta$ is reset to zero for the next
iteration (i.e.\ a steepest-descent step is taken), because the direction $\bt$
built on the old contact set is no longer valid for the new one.

\begin{remark}
  The overlap correction is absent from many informal descriptions of the
  Polonsky-Keer algorithm.
  For smooth surfaces (Hertz contact), convergence is only marginally
  affected by omitting it.  For rough surfaces, however, the contact set
  evolves non-monotonically: nodes can enter and leave $\calC$ repeatedly,
  and the overlap correction is essential for convergence.
\end{remark}

\subsection{Stopping criterion}

The complementarity error
\begin{equation}
  e_k = \frac{1}{P_{\mathrm{total}}\,g_{\mathrm{scale}}}
        \sum_{i=1}^N p_i\,|\tilde g_i|
  \label{eq:stopping}
\end{equation}
measures how far the current iterate deviates from the complementarity
condition~\eqref{eq:complementarity}.
Here $P_{\mathrm{total}} = N\bar p$ is the total load and
$g_{\mathrm{scale}} = \max(\bg_0) - \min(\bg_0)$ is a gap scale.
Convergence is declared when $e_k < \varepsilon_{\mathrm{tol}}$.

\subsection{Complete algorithm}

\begin{algorithm}[H]
\caption{Polonsky--Keer Projected CG with PR+ beta (P-K PCG)}
\label{alg:pk}
\KwIn{%
  Matvec $\bp \mapsto \bS\bp$;\;
  initial gap $\bg_0 \in \R^N$;\;
  nominal pressure $\bar p > 0$;\;
  tolerance $\varepsilon$;\;
  \texttt{use\_pr} $\in \{\text{true, false}\}$
}
\KwOut{ContactResult: $\bp^*$, $\delta^*$, iterations, error}
$\bp \leftarrow \bar p\,\bone$;\quad
$\bt \leftarrow \bzero$;\quad
$\tilde{\bg}^{\mathrm{prev}} \leftarrow \bzero$;\quad
$G_{\mathrm{old}} \leftarrow 1$;\quad
$\delta_{\mathrm{cg}} \leftarrow 0$\;
\For{$k = 0, 1, \ldots, \texttt{max\_iter}$}{
  $\bu \leftarrow \bS\bp$;\quad $\bg \leftarrow \bu + \bg_0$\;\Comment*{first matvec}
  $\calC \leftarrow \{i : p_i > 0\}$;\quad
  $\delta \leftarrow |\calC|^{-1}\!\sum_{i\in\calC} g_i$;\quad
  $\tilde\bg \leftarrow \bg - \delta$\;\Comment*{rigid-body shift}
  $e \leftarrow (P_{\mathrm{total}}\,g_{\mathrm{scale}})^{-1}\sum_i p_i|\tilde g_i|$\;
  \If{$e < \varepsilon$}{\textbf{break} \Comment*{converged}}
  $G \leftarrow \sum_{i\in\calC}\tilde g_i^2$\;
  \eIf{\texttt{use\_pr}}{
    $G_{\mathrm{pr}} \leftarrow \sum_{i\in\calC}\tilde g_i(\tilde g_i - \tilde g_i^{\mathrm{prev}})$\;
    $\beta \leftarrow \delta_{\mathrm{cg}}\cdot\max(0,\;G_{\mathrm{pr}}/G_{\mathrm{old}})$\;
  }{
    $\beta \leftarrow \delta_{\mathrm{cg}}\cdot G / G_{\mathrm{old}}$\;
  }
  $t_i \leftarrow \tilde g_i + \beta\,t_i$\; for $i \in \calC$;\quad
  $t_i \leftarrow 0$ for $i \notin \calC$\;
  $\tilde{\bg}^{\mathrm{prev}} \leftarrow \tilde\bg$;\quad $G_{\mathrm{old}} \leftarrow G$\;
  $\br \leftarrow \bS\bt$\;\Comment*{second matvec}
  $\bar r_\calC \leftarrow |\calC|^{-1}\sum_{i\in\calC}r_i$\;
  $\tau \leftarrow \dfrac{\sum_{i\in\calC}\tilde g_i\,t_i}
                         {\sum_{i\in\calC}(r_i - \bar r_\calC)\,t_i}$\;\Comment*{exact step, eq.~\eqref{eq:tau}}
  $p_i \leftarrow \max(p_i - \tau t_i,\;0)$ for $i \in \calC$\;\Comment*{project}
  \textsf{overlap} $\leftarrow$ false\;
  \For{$i$ with $p_i = 0$ and $\tilde g_i < 0$}{
    $p_i \leftarrow -\tau\tilde g_i$;\quad \textsf{overlap} $\leftarrow$ true\;\Comment*{overlap correction, eq.~\eqref{eq:overlap}}
  }
  $\delta_{\mathrm{cg}} \leftarrow 0$ if \textsf{overlap} else $1$\;\Comment*{reset CG on set change}
  $\bp \leftarrow \bp \cdot N\bar p\,/\!\sum_j p_j$\;\Comment*{enforce load balance}
}
\end{algorithm}

% ---------------------------------------------------------
\section{GPCG: Gradient Projected Conjugate Gradient}
\label{sec:gpcg}

The Gradient Projected Conjugate Gradient method of Mor\'e and
Toraldo~\cite{MoreToraldo1991} solves the general bound-constrained
QP~\eqref{eq:bqp} via an outer loop alternating between two phases:
(1) a gradient-projection phase that identifies active constraints, and
(2) a CG phase that minimises over the current free subspace.

\subsection{Phase 1: gradient projection}

Starting from $\by_0 = \bx_k$, the phase performs projected steepest-descent steps:
\begin{align}
  \alpha_{\mathrm{gp}} &= \frac{\norm{\nabla f(\by)}^2}{\nabla f(\by)^\tr \bA\,\nabla f(\by)},
  \label{eq:gp-step}\\[4pt]
  \by_{\mathrm{new}}   &= \Proj[\by - \alpha_{\mathrm{gp}}\nabla f(\by)],
\end{align}
until the free set $\calF(\by_{\mathrm{new}}) = \calF(\by)$ stabilises.
This identifies, in finite steps, the correct active set near the current iterate.

\subsection{Phase 2: CG on the free subspace}

Let $\bZ_k \in \R^{n \times |\calF_k|}$ be the restriction operator (columns are
the standard basis vectors $\be_i$ for $i \in \calF_k$).
The reduced problem is
\begin{equation}
  \min_{\bw \in \R^{|\calF_k|}}\; \tilde f(\bw)
  = \tfrac{1}{2}\bw^\tr(\bZ_k^\tr\bA\bZ_k)\bw
    + (\bZ_k^\tr\nabla f(\bx_k))^\tr\bw,
  \label{eq:reduced}
\end{equation}
solved approximately by CG, with the termination criterion of
Mor\'e and Toraldo: CG stops when the improvement in $\tilde f$
falls below $\eta_1$ times the largest previous improvement.
The step is then taken as
$\bx_{k+1} = \Proj(\bx_k + \alpha\bZ_k\bd)$
with $\alpha$ satisfying the Armijo condition (polynomial interpolation
backtracking, Section~\ref{sec:linesearch}).

\subsection{Convergence}

\begin{theorem}[Mor\'e \& Toraldo 1991]
  If $\bA$ is SPD and the active-set identification is finite, then GPCG
  converges to the unique solution $\bx^*$ of~\eqref{eq:bqp}.
  Moreover, once the optimal active set is identified, convergence is
  superlinear (q-superlinear for strongly convex $f$).
\end{theorem}

\subsection{Comparison with Polonsky--Keer}

\begin{table}[ht]
\centering
\caption{Comparison of the P-K PCG (Algorithm~\ref{alg:pk}) and GPCG algorithms.}
\label{tab:comparison}
\begin{tabular}{@{}lll@{}}
\toprule
Property & P-K PCG (Algorithm~\ref{alg:pk}) & GPCG~\cite{MoreToraldo1991} \\
\midrule
Loop structure       & Single loop               & Two-phase outer loop \\
Equality constraint  & \checkmark\; rigid-body shift + rescale & $\times$\; not handled \\
Overlap correction   & \checkmark\; essential for rough gaps & $\times$\; no equivalent \\
Active-set update    & Implicit (projection)     & Explicit ($\bZ_k$ matrix) \\
Line search          & Exact, 1 extra matvec     & Armijo backtrack, variable \\
Matvecs / outer iter & 2                         & $\geq 2$ \\
Convergence theory   & Heuristic $\beta$ reset   & Superlinear (strongly convex) \\
\bottomrule
\end{tabular}
\end{table}

\begin{remark}
  For the contact QP~\eqref{eq:contact-qp}, the mean-pressure equality
  constraint and the overlap correction make GPCG inapplicable without
  non-trivial modification.  The P-K algorithm, designed specifically for
  this problem, achieves 24--46 iterations in practice (Table~\ref{tab:iter}),
  which is entirely satisfactory given the $O(N\log N)$ cost of each matvec.
\end{remark}

% ---------------------------------------------------------
\section{Numerical Results}
\label{sec:numerics}

\subsection{Setup}

All runs use the Hertz contact gap
$g_0(x,y) = (x^2 + y^2)/(2R)$ with $R = 1$, domain size $L = 1$,
$E^* = 1$, and nominal pressure $\bar p = 0.05$.
The influence matrix is assembled as an H-matrix
($\eta = 2.0$, $\varepsilon_{\mathrm{ACA}} = 10^{-6}$, leaf size 32)
using the OpenMP-parallel ACA implementation.
All timings are single-run wall-clock times on a 20-core workstation.

\subsection{FR vs PR+ iteration counts}

\begin{table}[ht]
\centering
\caption{Comparison of Fletcher--Reeves (FR) and Polak--Ribi\`ere+ (PR+)
         $\beta$ formulas in Algorithm~\ref{alg:pk} at tolerance $\varepsilon = 10^{-8}$.}
\label{tab:iter}
\begin{tabular}{@{}rrrrrl@{}}
\toprule
$N_s$ & $N = N_s^2$ & FR iters & PR+ iters & $\Delta$ iters & FR solve / PR+ solve \\
\midrule
 64  &     4\,096 & 28 & 24 & $-14\%$ & 118\,ms / 99\,ms \\
128  &    16\,384 & 36 & 35 & $-3\%$  & 1024\,ms / 945\,ms \\
256  &    65\,536 & 46 & 44 & $-4\%$  & 3716\,ms / 3209\,ms \\
\bottomrule
\end{tabular}
\end{table}

PR+ consistently reduces the iteration count, with the largest benefit at
small grids where active-set changes are more frequent relative to the
total number of iterations.

\subsection{H-matrix scaling and large-grid results}

\begin{table}[ht]
\centering
\caption{H-matrix assembly time, memory, and solve time at $\varepsilon = 10^{-6}$
         (FR/PR+ iters with Algorithm~\ref{alg:pk}).}
\label{tab:scaling}
\begin{tabular}{@{}rrrrrlll@{}}
\toprule
$N_s$ & $N$ & Assembly & Memory & Compression & Avg.\ rank & FR iters & PR+ iters \\
\midrule
  64 &      4\,096 &    10\,ms &      40\,MiB & 0.31$\times$ & 11 & 28 & 24 \\
 128 &     16\,384 &    82\,ms &     $\sim$\!200\,MiB & --- & 11 & 36 & 35 \\
 256 &     65\,536 &   457\,ms &   $\sim$\!800\,MiB & --- & 11 & 46 & 44 \\
 512 &   262\,144 &   9.3\,s  &   6\,483\,MiB & 0.012$\times$ & 11 & 28 & 27 \\
\bottomrule
\end{tabular}
\end{table}

\noindent
Key observations:
\begin{itemize}
  \item The average ACA rank remains $\approx 11$ across all grid sizes,
        confirming the $O(N\log N)$ H-matrix complexity.
  \item The compression ratio improves dramatically with $N$: at $N_s = 512$
        the H-matrix uses 6.5\,GiB vs.\ 524\,GiB for a dense matrix.
  \item At $N_s = 512$ (262\,144 unknowns) the solver converges in $\approx 28$
        iterations; the solve time is dominated by the 9.3\,s assembly rather
        than the CG iterations.
  \item $N_s = 1024$ ($\approx 10^6$ unknowns) requires an estimated 26\,GiB
        for the H-matrix and is out of scope on the current workstation (12\,GiB RAM);
        it would require MPI or out-of-core ACA.
\end{itemize}

% ---------------------------------------------------------
\begin{thebibliography}{99}

\bibitem{AlBaali1985}
M.~Al-Baali.
\newblock Descent property and global convergence of the {Fletcher--Reeves}
  method with inexact line search.
\newblock \emph{IMA Journal of Numerical Analysis}, 5(1):121--124, 1985.

\bibitem{DaiYuan1999}
Y.-H.~Dai and Y.~Yuan.
\newblock A nonlinear conjugate gradient method with a strong global convergence property.
\newblock \emph{SIAM Journal on Optimization}, 10(1):177--182, 1999.

\bibitem{FletchReeves1964}
R.~Fletcher and C.~M.~Reeves.
\newblock Function minimization by conjugate gradients.
\newblock \emph{The Computer Journal}, 7(2):149--154, 1964.

\bibitem{GilNoc1992}
J.-C.~Gilbert and J.~Nocedal.
\newblock Global convergence properties of conjugate gradient methods for optimization.
\newblock \emph{SIAM Journal on Optimization}, 2(1):21--42, 1992.

\bibitem{HestStief1952}
M.~R.~Hestenes and E.~Stiefel.
\newblock Methods of conjugate gradients for solving linear systems.
\newblock \emph{Journal of Research of the National Bureau of Standards},
  49(6):409--436, 1952.

\bibitem{Love1929}
A.~E.~H.~Love.
\newblock The stress produced in a semi-infinite solid by pressure on part of the boundary.
\newblock \emph{Philosophical Transactions of the Royal Society A},
  228:377--420, 1929.

\bibitem{MoreToraldo1991}
J.~J.~Mor\'e and G.~Toraldo.
\newblock On the solution of large quadratic programming problems with bound constraints.
\newblock \emph{SIAM Journal on Optimization}, 1(1):93--113, 1991.

\bibitem{NocWri2006}
J.~Nocedal and S.~J.~Wright.
\newblock \emph{Numerical Optimization}.
\newblock Springer, 2nd edition, 2006.

\bibitem{PolKeer1999}
I.~A.~Polonsky and L.~M.~Keer.
\newblock A numerical method for solving rough contact problems based on the
  multi-level multi-summation and conjugate gradient techniques.
\newblock \emph{Wear}, 231(2):206--219, 1999.

\bibitem{PolakRib1969}
E.~Polak and G.~Ribi\`ere.
\newblock Note sur la convergence de m\'ethodes de directions conjugu\'ees.
\newblock \emph{Revue Fran\c{c}aise d'Informatique et de Recherche Op\'erationnelle},
  3(16):35--43, 1969.

\bibitem{Zoutendijk1970}
G.~Zoutendijk.
\newblock Nonlinear programming, computational methods.
\newblock In J.~Abadie, editor, \emph{Integer and Nonlinear Programming},
  pages 37--86. North-Holland, Amsterdam, 1970.

\end{thebibliography}

\end{document}
